\documentclass[12pt,a4paper]{article}
\usepackage[spanish]{babel}
\usepackage[utf8]{inputenc}
\usepackage[T1]{fontenc}
\usepackage{lmodern}
\usepackage{geometry}
\usepackage{setspace}
\usepackage{longtable}
\usepackage{booktabs}
\usepackage{array}
\usepackage{hyperref}
\usepackage{enumitem}

\geometry{margin=2.5cm}
\setstretch{1.2}
\hypersetup{colorlinks=true, linkcolor=black, urlcolor=blue, citecolor=black}

\title{Levantamiento de Documentacion del Proyecto\\\large Gestion Scrum para Plataforma Web de Mascotas 3D}
\author{Equipo del Proyecto}
\date{Mayo 2026}

\begin{document}

\maketitle
\tableofcontents
\newpage

\section{Introduccion}
Este documento presenta el levantamiento de documentacion del proyecto \textbf{Plataforma Web de Mascotas 3D} bajo el marco \textbf{Scrum}. El objetivo es definir de manera formal la organizacion del trabajo, los artefactos, los eventos, los roles y la trazabilidad del producto para el desarrollo incremental del \textit{MVP}.

El contenido se fundamenta en fuentes oficiales y reconocidas de Scrum, especialmente la \textit{Scrum Guide 2020} \cite{schwaber2020}, complementada por literatura especializada \cite{verheyen2020,pichler2010}.

\section{Contexto del Proyecto}
\subsection{Vision del producto}
Construir una plataforma web donde una persona usuaria pueda:
\begin{itemize}[leftmargin=*]
    \item Registrarse e iniciar sesion.
    \item Registrar informacion de sus mascotas (nombre, edad, fotos y datos relevantes).
    \item Asociar y visualizar modelos 3D base de animales.
\end{itemize}

El alcance inmediato corresponde al \textbf{MVP}: autenticacion y registro de mascotas. La personalizacion avanzada de modelos 3D se planifica como evolutivo posterior.

\subsection{Objetivo del levantamiento}
\begin{itemize}[leftmargin=*]
    \item Definir la gestion del proyecto con Scrum de forma auditable.
    \item Establecer artefactos y eventos con criterios claros.
    \item Alinear el trabajo tecnico con objetivos de producto y entrega incremental.
\end{itemize}

\section{Fundamento Teorico de Scrum}
\subsection{Definicion}
Scrum es un marco ligero para desarrollar, entregar y sostener productos complejos, apoyado en empirismo y \textit{lean thinking} \cite{schwaber2020}. El control empirico se basa en tres pilares:
\begin{itemize}[leftmargin=*]
    \item \textbf{Transparencia}
    \item \textbf{Inspeccion}
    \item \textbf{Adaptacion}
\end{itemize}

\subsection{Valores Scrum}
Segun la Scrum Guide, Scrum se sostiene en cinco valores \cite{schwaber2020}:
\begin{itemize}[leftmargin=*]
    \item Compromiso
    \item Foco
    \item Apertura
    \item Respeto
    \item Coraje
\end{itemize}

\section{Roles Scrum para el Proyecto}
\subsection{Product Owner}
Responsable de maximizar valor del producto y gestionar el \textit{Product Backlog}. Define prioridades y criterios de aceptacion.

\subsection{Scrum Master}
Facilita eventos Scrum, remueve impedimentos y asegura la correcta aplicacion del marco.

\subsection{Developers}
Equipo multidisciplinario que transforma elementos del \textit{Product Backlog} en incrementos potencialmente liberables durante cada Sprint.

\section{Eventos Scrum y Cadencia}
\begin{longtable}{>{\raggedright\arraybackslash}p{3.2cm} >{\raggedright\arraybackslash}p{3.0cm} >{\raggedright\arraybackslash}p{7.0cm}}
\toprule
\textbf{Evento} & \textbf{Duracion} & \textbf{Objetivo en el proyecto} \\
\midrule
Sprint & 2 semanas & Entregar incremento funcional del MVP. \\
Sprint Planning & 2 horas & Definir Sprint Goal y seleccionar PBIs. \\
Daily Scrum & 15 minutos & Sincronizar avance y detectar bloqueos. \\
Sprint Review & 1.5 horas & Inspeccionar incremento con interesados y obtener feedback. \\
Sprint Retrospective & 1 hora & Mejorar proceso y acuerdos de trabajo. \\
\bottomrule
\end{longtable}

\section{Artefactos Scrum del Proyecto}
\subsection{Product Backlog}
Lista ordenada y dinamica de trabajo pendiente para el producto \cite{schwaber2020}. Para este proyecto, las epicas principales son:
\begin{itemize}[leftmargin=*]
    \item Autenticacion y gestion de usuarios.
    \item Registro y gestion de mascotas.
    \item Visualizacion de modelos 3D base.
    \item Seguridad y soporte tecnico transversal.
\end{itemize}

\subsection{Sprint Backlog}
Conjunto de PBIs seleccionados para el Sprint, mas el plan de entrega. Incluye tareas tecnicas, pruebas y validaciones.

\subsection{Incremento}
Resultado integrado al final del Sprint que cumple la \textit{Definition of Done} y aporta valor utilizable.

\subsection{Commitments oficiales por artefacto}
\begin{itemize}[leftmargin=*]
    \item Product Backlog $\rightarrow$ \textbf{Product Goal}
    \item Sprint Backlog $\rightarrow$ \textbf{Sprint Goal}
    \item Incremento $\rightarrow$ \textbf{Definition of Done}
\end{itemize}

\section{Aplicacion Scrum al MVP}
\subsection{Product Goal}
Entregar un MVP funcional en el que una persona usuaria pueda crear cuenta, iniciar sesion y registrar mascotas con informacion basica y fotos, incluyendo visualizacion de modelos 3D base asociables.

\subsection{Sprint Goal propuesto (Sprint 1)}
Implementar autenticacion completa (registro/login) y alta de mascotas con persistencia en base de datos.

\subsection{PBIs priorizados para Sprint 1}
\begin{enumerate}[leftmargin=*]
    \item Registro de usuario con validacion y cifrado de contrasena.
    \item Inicio de sesion con token JWT.
    \item Creacion de mascota (nombre, edad, tipo).
    \item Listado de mascotas por usuario autenticado.
    \item Validaciones de seguridad y errores de entrada.
\end{enumerate}

\subsection{Definition of Done (DoD)}
Un item se considera \textit{Done} cuando cumple:
\begin{itemize}[leftmargin=*]
    \item Codigo integrado sin errores de compilacion.
    \item Pruebas unitarias minimas acordadas por el equipo.
    \item Endpoint validado funcionalmente.
    \item Criterios de aceptacion cumplidos.
    \item Sin vulnerabilidades evidentes de autenticacion/autorizacion.
\end{itemize}

\section{Trazabilidad de Requisitos Scrum}
\begin{longtable}{>{\raggedright\arraybackslash}p{1.6cm} >{\raggedright\arraybackslash}p{5.8cm} >{\raggedright\arraybackslash}p{2.7cm} >{\raggedright\arraybackslash}p{3.1cm}}
\toprule
\textbf{ID} & \textbf{PBI} & \textbf{Prioridad} & \textbf{Sprint objetivo} \\
\midrule
PBI-01 & Registro de usuarios & Alta & Sprint 1 \\
PBI-02 & Login con JWT & Alta & Sprint 1 \\
PBI-03 & Registro de mascotas & Alta & Sprint 1 \\
PBI-04 & Carga de fotos de mascota & Media & Sprint 2 \\
PBI-05 & Visualizacion modelo 3D base & Media & Sprint 2 \\
PBI-06 & Personalizacion 3D avanzada & Baja & Post-MVP \\
\bottomrule
\end{longtable}

\section{Riesgos y Controles en Scrum}
\subsection{Riesgos}
\begin{itemize}[leftmargin=*]
    \item Sobrecarga de alcance en Sprint (\textit{scope creep}).
    \item Dependencias tecnicas en carga/render 3D.
    \item Riesgos de seguridad en autenticacion.
\end{itemize}

\subsection{Controles}
\begin{itemize}[leftmargin=*]
    \item Priorizacion estricta por valor de negocio.
    \item Revisiones incrementales con demo funcional.
    \item Retrospectivas con acciones concretas de mejora.
\end{itemize}

\section{Indicadores de Seguimiento}
\begin{itemize}[leftmargin=*]
    \item Velocidad del equipo (story points por Sprint).
    \item Porcentaje de PBIs completados por Sprint.
    \item Defectos encontrados en Sprint Review.
    \item Tiempo de ciclo por historia.
\end{itemize}

\section{Conclusiones}
La adopcion de Scrum permite estructurar el desarrollo del proyecto con entregas incrementales, inspeccion continua y adaptacion basada en evidencia. Para este proyecto, Scrum habilita concentrar el esfuerzo en un MVP realista (autenticacion y registro de mascotas), dejando la personalizacion 3D avanzada como expansion posterior de forma controlada.

\newpage
\begin{thebibliography}{9}

\bibitem{schwaber2020}
Ken Schwaber and Jeff Sutherland.
\textit{The Scrum Guide: The Definitive Guide to Scrum: The Rules of the Game}.
Scrum.org and ScrumInc., 2020.
Available at: \url{https://www.scrum.org/resources/scrum-guide}

\bibitem{verheyen2020}
Gunther Verheyen.
\textit{Scrum -- A Pocket Guide: A Smart Travel Companion}.
Van Haren Publishing, 3rd edition, 2020.

\bibitem{pichler2010}
Roman Pichler.
\textit{Agile Product Management with Scrum: Creating Products that Customers Love}.
Addison-Wesley, 2010.

\end{thebibliography}

\end{document}
